\section{模型评价与推广}

\subsection{模型优点}

本文模型具有三点优势。第一，四问共享同一组温度、水分、环境边界和通量符号约定，模型从固定域状态场自然递进到移动域全域事件，避免重复定义变量。第二，有限体积离散直接守恒径向界面通量，表面半网格重构保留了真实表面状态，能够同时满足题目指定位置输出和全域达标判定。第三，问题四将附件2的半径轨迹作为几何输入，并以材料坐标处理收缩，固定物理位置出域后留空、表面列随实时半径更新，输出口径与几何模型一致。

\subsection{模型局限}

模型只描述圆柱主体的径向过程，不能刻画端面附近的轴向效应；长度不变和材料位置同比例缩放是由数据条件出发的补充假设。模型未建立含水率与尺寸收缩之间的本构关系，而将附件2给出的 $R(t)$ 作为外部几何输入；能量方程未显式考虑蒸发潜热，模型也未加入吸附等温线及 Soret/Dufour 等热质交叉机制。环境末值延拓同样会影响长时段预测。经验密度关系只进入热物性，不被用来反推几何尺寸。因此，本文结果应解释为给定输入、边界与简化条件下的数值预测，而不是对所有实际药材和工艺的普适结论。

\subsection{模型推广}

若获得端面温度或水分数据，可将一维径向模型扩展为轴向—径向二维模型，并保留当前的守恒离散和全域事件判定。若有独立的收缩方向数据，可把长度变化和径向变化分别写入材料坐标映射；若能测得潜热、吸附等温线和收缩速率，则可在能量方程和质量方程中加入相应耦合项。对于新的环境边界，只需替换经过审查的时间序列并重新进行网格、守恒和边界根检验，不能直接沿用本题的临界时间。
